\section{Appendix}
\label{sec:appendix}

\noindent This appendix retains only compact audit material that supports the accepted empirical and DTSE chapters. Core chapter argumentation, scenario framing, and methodology exposition remain in the main body.

\subsection{DTSE Scenario-Formula Map}
\label{app:dtse_scenario_formula_map}
Table~\ref{tab:app_dtse_scenario_formula_map} gives a compact reader-facing summary of why each scenario is included, how it enters the evaluator, and which fixed settings matter most for interpretation.

\begin{table}[H]
    \centering
    \small
    \renewcommand{\arraystretch}{1.15}
    \setlength{\tabcolsep}{4pt}
    \begin{tabularx}{\textwidth}{@{}>{\raggedright\arraybackslash}p{0.18\textwidth}>{\raggedright\arraybackslash}p{0.20\textwidth}>{\raggedright\arraybackslash}p{0.27\textwidth}X@{}}
        \toprule
        \textbf{Scenario} & \textbf{Why Included} & \textbf{Rule Summary} & \textbf{Key Fixed Settings} \\
        \midrule
        Baseline Neutral &
        Establish the reference path for timing and magnitude comparisons. &
        Stable demand and background settings generate the neutral baseline against which all later deviations are measured. &
        No adverse event trigger; no signal threshold needs to be crossed for the baseline to exist as a reference path. \\
        \addlinespace
        Demand Contraction &
        Test reward--demand decoupling under weakening service use. &
        Demand decays through the scenario-specific contraction path under higher volatility. &
        Utilization becomes a material signal once it departs ten percent from the baseline path. \\
        \addlinespace
        Liquidity Shock &
        Test fast transmission from token-liquidity stress into provider-economics fields. &
        A discrete unlock-style sell event is introduced under limited market depth, producing a sudden price dislocation. &
        Event week: 20. Sell-pressure scalar: 35\%. Material price deviation threshold: 10\%. \\
        \addlinespace
        Competitive Yield Pressure &
        Isolate provider reallocation pressure without requiring demand collapse. &
        External yield pressure raises switching sensitivity, with provider types responding at different intensities under the same outside opportunity. &
        External-yield scalar: 1.5$\times$. Provider-count threshold: 5\%. Churn-increase threshold: 20\%. \\
        \addlinespace
        Provider Cost Inflation &
        Test Subsidy-Gap exposure under higher operating costs. &
        Provider cost baselines are shifted upward across the active provider set while all other scenario channels remain neutral. &
        Cost multiplier: 1.35$\times$. Profitability threshold: 10\%. Provider-count threshold: 5\%. \\
        \bottomrule
    \end{tabularx}
    \caption{Scenario logic and key fixed settings for the DTSE experiment set.}
    \label{tab:app_dtse_scenario_formula_map}
\end{table}

\subsection{DTSE Threshold Definitions}
\label{app:dtse_threshold_table}
Table~\ref{tab:app_dtse_threshold_registry_excerpt} summarizes the main fixed thresholds referenced in Sections~\ref{sec:dtse_methodology} and~\ref{sec:simulation_results}.

\begin{table}[H]
    \centering
    \footnotesize
    \renewcommand{\arraystretch}{1.15}
    \setlength{\tabcolsep}{4pt}
    \begin{tabularx}{\textwidth}{@{}>{\raggedright\arraybackslash}p{0.27\textwidth}>{\raggedright\arraybackslash}p{0.11\textwidth}>{\raggedright\arraybackslash}p{0.12\textwidth}X@{}}
        \toprule
        \textbf{Threshold or Setting} & \textbf{Value} & \textbf{Unit} & \textbf{Interpretive Role} \\
        \midrule
        Negative-provider-profit rule & 0 & weekly profit & Separates non-negative provider economics from loss-making periods in profitability-based readouts. \\
        \addlinespace
        Sustained-loss window & 3 & weeks & Defines how long negative profitability must persist before profitability-induced churn is treated as a Stage~2 pattern. \\
        \addlinespace
        Liquidity event timing & 20 & week index & Marks the discrete sell event used in the liquidity-shock scenario. \\
        \addlinespace
        Liquidity sell-pressure scalar & 35 & percent & Defines the size of the unlock-style sell event in the liquidity-shock scenario. \\
        \addlinespace
        Competitive-yield scalar & 1.5 & relative multiple & Sets the strength of the outside opportunity used in the competitive-yield scenario. \\
        \addlinespace
        Provider-cost multiplier & 1.35 & multiplier & Sets the exogenous increase in operating costs used in the cost-inflation scenario. \\
        \addlinespace
        Price-signal threshold & 10 & percent & Defines a material price deviation from baseline in Stage~1 detection. \\
        \addlinespace
        Utilization-signal threshold & 10 & percent & Defines a material utilization deviation from baseline in Stage~1 detection. \\
        \addlinespace
        Profitability-signal threshold & 10 & percent & Defines a material profitability deviation from baseline in Stage~1 detection. \\
        \addlinespace
        Provider-count threshold & 5 & percent & Defines a material active-provider deviation from baseline in Stage~1 detection. \\
        \addlinespace
        Churn-increase threshold & 20 & percent & Defines a material churn deviation from baseline in Stage~1 detection. \\
        \bottomrule
    \end{tabularx}
    \caption{Excerpt of the fixed thresholds used in the DTSE reporting logic.}
    \label{tab:app_dtse_threshold_registry_excerpt}
\end{table}

\subsection{Override Ledger}
\label{app:dtse_override_ledger}
The reported DTSE runs also include a small number of profile-specific and global overrides. Table~\ref{tab:app_dtse_override_ledger} summarizes the main ones in human-readable form.

\begin{table}[H]
    \centering
    \small
    \renewcommand{\arraystretch}{1.15}
    \setlength{\tabcolsep}{4pt}
    \begin{tabularx}{\textwidth}{@{}>{\raggedright\arraybackslash}p{0.24\textwidth}>{\raggedright\arraybackslash}p{0.18\textwidth}>{\raggedright\arraybackslash}p{0.16\textwidth}X@{}}
        \toprule
        \textbf{Override Family} & \textbf{Scope} & \textbf{Status} & \textbf{Purpose} \\
        \midrule
        Profile-specific churn thresholds & ONO, Helium, Geodnet & Modeled assumption & Allow continuation sensitivity to differ across mechanism profiles without treating those values as empirical estimates of live behavior. \\
        \addlinespace
        ONO maximum weekly churn cap & ONO & Modeled assumption & Prevents unrealistic discontinuities in weekly exit intensity while preserving participation stress under adverse conditions. \\
        \addlinespace
        Provider-type switching adjustment & Global competitive-yield rule & Scenario logic & Allows competitive-yield pressure to affect provider types differently under the same outside opportunity. \\
        \bottomrule
    \end{tabularx}
    \caption{Main profile-specific and global overrides used in the reported DTSE runs.}
    \label{tab:app_dtse_override_ledger}
\end{table}

\subsection{DTSE Run Manifest}
\label{app:dtse_run_manifest}
Table~\ref{tab:app_dtse_manifest_excerpt} records the key features of the frozen run manifest used for Sections~\ref{sec:simulation_results} and~\ref{sec:discussion}.

\begin{table}[H]
    \centering
    \small
    \renewcommand{\arraystretch}{1.15}
    \setlength{\tabcolsep}{5pt}
    \begin{tabularx}{\textwidth}{@{}>{\raggedright\arraybackslash}p{0.32\textwidth}X@{}}
        \toprule
        \textbf{Manifest Field} & \textbf{Recorded Value} \\
        \midrule
        Run timestamp (UTC) & 24 February 2026, 10:21 UTC \\
        Scenario-grid identifier & DTSE Chapter 7 grid, version 1 \\
        Deterministic seed policy & Master seed plus run index under a fixed seed schedule \\
        Horizon and repetitions & 52 weeks per run; 60 repeated runs per scenario--profile condition \\
        Mechanism profiles & ONO calibrated profile, Helium BME profile, Hivemapper profile, Grass profile, and Geodnet profile \\
        Scenario set & Baseline Neutral, Demand Contraction, Liquidity Shock, Competitive Yield Pressure, and Provider Cost Inflation \\
        \bottomrule
    \end{tabularx}
    \caption{Run-manifest excerpt for the frozen DTSE reporting set.}
    \label{tab:app_dtse_manifest_excerpt}
\end{table}

\paragraph{Empirical applicability note}
The empirical metric-applicability boundary remains in Section~\ref{sec:empirical_stress_layer}, where the thesis defines the \emph{N/R} rule, applies the comparator-level matrix, and states the Onocoy-specific observability limits directly in the main body rather than repeating them in appendix form.
